\section{Ashura}

These extracts from the longer work
\emph{Persian Traditions

in Spain\footnote{\url{https://www.gornahoor.net/?page\_id=11751}}} by \textbf{Michael McClain} were originally published in the
\emph{International Journal of Shi'i Studies}, Volume V, No. 1, 2007,

under the title ``Affinities Between Catholicism and Shi'ism''. This is part 4 of 4.

Mr. McClain concludes with the similarities between Persian Shi'ism, Russian Orthodoxy, and Spanish Catholicism.

← Part 3 The White Dove\footnote{\url{https://gornahoor.net/?p=8728}}

The parallels between Iranian Shi`ism, on the one hand, and Spanish Catholicism and to a slightly lesser degree Irish Catholicism, on the other, is a topic too vast to fully treat here. The resemblance between them is indeed uncanny. As we have said earlier, at least since the twelfth century, it has been noted repeatedly that Shi`a Islam has many special affinities with traditional Catholicism and Eastern Orthodox Christianity.

Spanish Catholicism has particular characteristics which give it a special rapport with Shi`ism above and beyond the general Shi`a-Catholic affinities. This is easily explained historically, as we have said throughout this chapter.

The Russian Orthodox Church also has particular characteristics which give it a special rapport with Shi`ism above and beyond the general Shi`a-Eastern Orthodox affinities. It is for this reason that I have devoted so much space in this book to Russia and Ukraine and to the Russian Orthodox Church.

It is also evident that, when speaking of Muslim-Christian dialogue, the ideal spokesman from the Christian side would be either a mystically oriented Spanish Catholic or a mystically oriented Russian Orthodox believer, while from the Muslim side the ideal spokesman would be a Sufi oriented Shi`a.

In a personal communication, my friend Seyyed Hossein Nasr told me: ``You are completely right in emphasizing the unique rapport between Shi`ism and Sufism on the one hand and certain elements of Spanish Catholicism and Russian Orthodoxy on the other.'' Since Spanish Catholicism and Russian Orthodoxy have particular characteristics which give them special affinities with Shi`ism, one may ask if they have special affinities among themselves. This is a potentially large topic which would require much research, so we will touch on it only briefly... .

At this point I can do no better than quotes the words of Sarah Hobson, speaking of her visit to Qum:

\begin{quotex}
I wanted to stay in Qum to participate further in this (Iranian Shi`a) religious life, and to learn more about Islam. I felt I had touched only the surface, even though that surface seemed misleadingly smooth, misleadingly simple. So many of the rules, the recommendations, the beliefs of Shi`a Islam seemed clear and easy to grasp. Yet I felt that below there was something far more complex, more intricate, of philosophies and theological arguments which were beyond my understanding. Perhaps, if I could have more time here (in Qum), I would manage to go deeper, to understand the mystical undertones of Islam, its esoterics, its thought patterns. Perhaps I would just understand, for I felt here in Qum that the religious leaders and teachers were closer to knowledge of truth than I had encountered in people elsewhere.\hyperref[_edn1]{ i }
\end{quotex}


My time in India convinced me that books are no substitute for a living Pir. I, too, should like to spend time in Qum, to learn of Shi`a Islam in depth, of its mysticism and esoterism. In Delhi I gave a talk on Muslim Spain at the tomb of Shaykh Nizamuddin. I remember my audience was composed mostly of very saintly looking men with white beards. Afterwards I was decorated with a leis of marigolds.

The journalist Robin Wright tells of her visit to Qum:

In stark contrast to the simplicity of (Ayatollah) Khomeini's barren home was a second religious landmark in Qum. The bustling new computer center at the Golpaygani Seminary (is) an impressively yellow brick building with blue tiled trim... .The computer center (in Qum) was the brainchild of Ayatollah `Ali Korani, a gentle cleric with a white beard and a white turban who had not even known how to type when he decided, about the time of Ayatollah Khomeini's death in 1989, that (Shi`a) Islam had to meet modernity in the form of computers. The idea grew out of his own rather unusual research. ``Since I was a child I have loved Imam Mahdi, our Twelfth Imam who disappeared,'' he (Ayatollah Korani) told me when I called on him at the computer center, where he and a young cleric were engrossed in a new software program. ``Some of us believe from reading Islamic sources that he (Imam Mahdi) will come back at the same time as Jesus. The coming of the Mahdi is certain because the Prophet Muhammad said it would happen. He will come first and go to Quds (Jerusalem). Then Jesus will appear from Heaven and they will pray in Quds together. And then there will be unity between Christians and Muslims,'' Korani explained.''\hyperref[_edn2]{ ii }

In Spain, mysticism is the lifeblood of religion, whether Muslim or Christian, and no doubt it was the lifeblood of Druidism in Spain in pre-Roman, pre-Christian times. This fact is of crucial importance, because nearly all the population of al-Andalus was composed either of Spanish Catholics, i.e., the Mozarabs, or of descendants of Spanish Catholics converted to Islam. In this last category must be included Ibn Hazm of Cordoba, Ibn `Arabi al-Mursi, Ibn `Abbad of Ronda, and the poets Ibn Quzman (pronounced ``Guzman'' in Spain) and Abu Bakr Ibn al-Qutiyya. The list could go on and on and on. In fact, nearly all the great figures of the history of al-Andalus as well as nearly all the general population either were Mozarabs or were descendants of Mozarabs converted to Islam. Even the Caliphs of Cordoba are included in this category, because far more Iberian, Celtic and Visigothic blood flowed in their veins than did Arab blood. Did the Caliphs of Cordoba themselves feel a twinge of ``Shi`a tendencies,'' ironically helping to provoke the fear of Fatimid subversion?

There is another fact which seems to indicate Shi`ism, open or clandestine, or at the very least strong Shi`a tendencies among the population of al-Andalus and, perhaps, even in the case of the great Caliph `Abd al-Rahman III. I refer to the prevalence of the feminine name \emph{Zahra,} usually pronounced ``Zahara'' or ``Zahira'' among Hispano-Muslims. The title ``Zahra'' for Fatima, daughter of the Prophet Muhammad, is extremely common among Shi`as. I would not go so far as to say that said title is unknown or never used among Sunnis, though I have never heard it said by a Sunni nor have I read it in Sunni literature. At the very least, the title ``Zahra'' for Fatima, daughter of the Prophet Muhammad, is very rarely or never used by Sunnis. Usually in the form of the transcriptions of Hispano-Muslim vocalizations ``Zahara'' or ``Zahira,'' the title ``Zahra'' is found to this day as a feminine personal name and in place names in Southern Spain. We have already mentioned the caliphal city near Cordoba, ``Medina al-Zahara,'' which is ``Madinat al-Zahra'' in correct, classical Arabic, built by the Caliph `Abd al-Rahman III. I also recall visiting the picturesque present day village called ``Zahara,'' not far from Ronda.

Over the main gate of the Caliphal residence in Medina al-Zahara was a statue of a woman, said to have been that of a ``favorite'' of the Caliph named ``Zahara.'' (156) Did this statue really represent a ``favorite'' of the Caliph or did it represent Fatima Zahra? Should this be the case, it would indicate that the Caliph `Abd ar-Rahman III was himself a clandestine Shi`a. This is all the more credible when we recall that said residence was within the caliphal city called ``Madinat al-Zahara.'' Would the great Caliph `Abd al-Rahman III really have named a city, and perhaps some villages and small towns, after a ``favorite,'' and adorned the main gate to the caliphal residence in Madinat al-Zahara with a statue of said favorite?

It is also said that in Cordoba in Muslim times there was a statue of the Virgin Mary above that gate of the city called ``Bab al-Qantara,'' which led to the bridge over the Guadalquivir, and another such statue over the principal gate of Pechina, near Almeria.\hyperref[_edn3]{ iii } We have already mentioned the many connections between the Virgin Mary and Fatima Zahra. Did these statues in cities of Muslim Spain really represent the Virgin Mary or did they represent Fatima Zahra, or, perhaps, had the two become fused in the popular imagination? I have never heard of any statues of Muslim holy men in Muslim Spain.

As is evident from what we have said above, al-Andalus was a land of Sufism, being second only to Persia and perhaps Muslim India in this respect. Sufism saturated the very air of al-Andalus. Indeed it was Sufism which won the bulk of the population of al-Andalus to Islam... .

It has been noted that some Shi`a mullahs have been violently anti-Sufi, particularly during the Safavid period, and this apparent fact has been used to ``prove'' that there is no intrinsic relation between Shi`ism and Sufism, though of course many Shi`as are Sufis and many Sufis are Shi`as. However, by this standard Ibn `Arabi al-Mursi was also anti-Sufi, since he strongly condemned some of the beliefs and practices of certain Sufis.\hyperref[_edn4]{ iv } This is a bit like saying that the Pope is anti-Catholic because he disapproves of some of the beliefs and practices of certain people who call themselves Catholics.

The truth is that, whether called \emph{Batin}, \emph{Hikmat} or
\emph{Irfan}, Sufism is an absolutely vital, organic, and indispensable

part of Shi`a doctrine, theology, philosophy and spirituality.\hyperref[_edn5]{ v } The same cannot be said of Sunni Islam, though, as we have said, it is most certainly true that there are Sunnis as well as Shi`as who are Sufis. To summarize, a Sunni Muslim may be a complete non-Sufi, having no connection whatever with Sufism in any way, shape or form. However, this is not possible for a Shi`a, who by virtue of being a Shi`a is a Sufi in some degree or another, though this degree varies greatly in individual cases. Conversely, a Sufi, though he may confessionally be a Sunni, by the fact of being a Sufi has certain ``Shi`a tendencies.''... We have already mentioned the parallels between Holy Week in Spain and the Shi`a Ashura. During my years in Granada I heard of celebrations of Good Friday which included auto flagellation and representations of the Crucifixion which were a bit too realistic, leading to serious injury or even death. I was not able to confirm the veracity of these tales. Apparently all this occurred in remote and inaccessible areas of the provinces of Granada, Malaga and Almeria, always areas with a strong Morisco background. Shi`ism would have certainly had different nuances among the Mudejares and most especially the Moriscos than among the Hispano-Muslims. The fact of living in an overwhelmingly Catholic environment would have certainly strengthened any Shi`a tendencies among the Moriscos for reasons given above, though in an unconscious manner and with no intention on anyone's part. One can readily imagine a Morisco or even a Mudejar participating in Good Friday celebrations, saying to himself: ``I do not believe that the Prophet Isa (Jesus Christ), on Whom be peace, was crucified, but Imam Hussayn, on Whom be peace, was most certainly cruelly martyred, and both the Prophet Isa and Imam Hussayn are alive and awaited, to return on the Day of Judgment, so in my heart I celebrate the martyrdom of Imam Hussayn, and with all sincerity I have great reverence for the Prophet Isa.''

In remote Morisco villages where even the parish priest might have been a Morisco, one can readily imagine practices typical of Ashura being used during the Good Friday celebrations. It may be that some of the typically Spanish customs connected with Good Friday and Holy Saturday celebrations may have been inspired by Shi`a influence by way of Mudejares and Moriscos. The history of the Franciscans in what is now the southwestern United States is an epic of heroism and sanctity, with an abundant harvest of martyrs. During the colonial period the Church in New Mexico was under the jurisdiction of the Bishop of Durango in Mexico. In practice New Mexico was a Franciscan province, because between Durango and New Mexico was 1500 miles of very rugged terrain full of Apaches. The Franciscans were in charge of the missions among the Amerindians as well as the spiritual well-being of the white colonists. A large proportion of the white males of New Mexico were Franciscan Lay Brothers.

The Mexican Period is not favorably remembered in New Mexico. So unpopular was the Mexican government that there were two uprisings against it in New Mexico. In 1828 the Mexican Government expelled all clergy of Spanish birth and secularized the mission lands. This meant that New Mexico was left almost without priests. While the Amerindians in many cases returned to pagan practices, something surprising occurred among the white population. Quite suddenly there appeared the brotherhoods of the Brothers of the Light and the Brothers of the Blood. During Holy Week these brotherhoods did things which the Franciscans would never have permitted, such as auto flagellation until the blood flowed and severely realistic representations of the Crucifixion. The popular name of these brotherhoods was, and is, \emph{Los Penitentes.}

In 1942 the Penitentes promised to cease and desist from the most extreme of their practices, and were accepted by the Church. Today in New Mexico about 2,500 men belong to the Penitente Brotherhoods. During Holy Week the Penitentes dress in white shorts and black hoods. To the sound of a \emph{pitero}, or small flute, they bear heavy crosses to the tops of the hills. They lie down on thorns and cactus spines. They flagellate themselves with yucca whips, at times lacerating the flesh until reaching the bone. They cut themselves with knives of obsidian or volcanic glass. On Good Friday the Crucifixion is represented, with a volunteer tied to a cross. At times the volunteer loses consciousness. There are stories of men being nailed to crosses and even dying as a result of a representation a bit too realistic.\hyperref[_edn6]{ vi }

Julio Puyol says that auto flagellation became part of Holy Week celebrations in Spain only in the first third of the Sixteenth Century. In 1565-66 the Valencia Provincial Council mentions the Thursday and Good Friday abuses (i.e., auto flagellation) of the penitentes.\hyperref[_edn7]{ vii } This is extremely interesting, because until a series of decrees around the year 1500, the Mudejares of Spain were openly Muslims. In other words, the above indicates that auto flagellation and similar Holy Week practices became common in Spain precisely at the time when the former Mudejares, now ``Moriscos'' or at least nominal Christians, would have become participants in said Holy Week ceremonies. Also note that Valencia, together with Granada, was the region with the highest density of Morisco population. All this points to a Muslim origin for many Spanish Holy Week customs, including auto flagellation. Among Muslims, these practices are typical only of Shi`as.

The Iranian writer Roy Mottahedeh has also noted the close parallels between the Holy Week celebrations of the Penitentes of New Mexico and the Ashura celebrations of the Shi`as of Iran. He says:

Flagellation survives in Spain and in many parts of the Hispanic world. It survives, in fact, in the United States in New Mexico, where, in spite of a century of horrified disapproval of Protestants and non-Hispanic Catholics, the Brotherhoods of Penitentes commemorate the Passion of Jesus by flagellation, by the carrying of heavy wooden crosses, and many other forms of discipline, physical and spiritual. The resemblance of the form of Penitente religiosity to (Iranian) Shiah practices extends to tableaux from the life of Jesus and even to the drama of a simulated crucifixion.

The resemblance in psychological content is even more striking: both the (Iranian) Shiah and the New Mexican Penitentes are using violation of physical self-integrity as a means to enter an altered state of awareness in which ordinary restraints of prudence are removed and the Penitente loses not only his sense of self-protection but also his sense of separateness. By sharing his ``discipline'' the Penitente has broken the boundary between himself and his fellow Penitentes and even --- to some extent --- between himself and the spiritual model he seeks to imitate; as the Penitentes say, the ``Brothers of Blood'' become ``Brothers of Light.'' At bottom, both are forms of folk mysticism.\hyperref[_edn8]{ viii }

The last paragraph hits the nail on the head. What the Penitentes of New Mexico really represent is Shi`ite religiosity expressed in Christian terms. Note also that Penitente religiosity suddenly appeared full-blown at the moment in which the Church in New Mexico was passing through a very difficult period, and returned to clandestinity or semi-clandestinity when an energetic group of French and Spanish priests under Bishop Lamy restored the Church in New Mexico. This indicates a long history and a long experience of clandestinity and \emph{taqiyya}, or dissimulation. In Muslim Spain, the Shi`ites usually had to practice at least a semi-clandestinity, as in Persia and other Middle Eastern countries under dynasties of Turkish origin. Under the kings of Castile, the Church suppressed any attempts to use the more extreme practices of Ashura in the Holy Week celebrations, except, perhaps, in some very remote and inaccessible places, but through clandestinity something survived. With the passage of so many generations between 1492 (the year that the last Nazirid king surrendered Granada to Fernando of Aragon and Isabel of Castile) and 1828 the Imam Hussayn was forgotten and Jesus Christ took his place, but this change was purely nominal, the religiosity itself remaining the same.

\hyperref[_ednref1]{ i } Sarah

Hobson, \emph{Through Iran in Disguise} (Chicago, 1982), p. 68.

\hyperref[_ednref2]{ ii } Robin

Wright, \emph{The Last Great Revolution} (New York, 2000), p. 236.

\hyperref[_ednref3]{ iii } Henri

Peres, \emph{Esplendor de al-Andalus} (Madrid, 1983), p. 333.

\hyperref[_ednref4]{ iv } Ibn

Arabi, \emph{Risalat al-Quds}, translated as \emph{Vidas de Santones Andaluces} by Miguel Asin Palacios (Madrid, 1981), pp. 121-122.

\hyperref[_ednref5]{ v } See

\emph{The Heritage of Sufism, Volume I: Classical Persian Sufism from

its Origins to Rumi (700-1300)}, edited by Leonard Lewisohn (Oxford, 1999), essay by John Cooper, ``Rumi and Hikmat: Towards a Reading

of Sabziwari's Commentary on the Mathnawi,'' pp. 409-435. \emph{The Heritage of Sufism, Volume II: The Legacy of Medieval Persian Sufism (1150-1500)}, edited by Leonard Lewisohn (Oxford, 1999), essay by Leonard Lewisohn, ``Overview: Iranian Islam and Persianate Sufism,'' p. 19. \emph{The Heritgae of Sufism, Volume III: Late Classical Persianate Sufism (1501-1750)} essays by Seyyed Hossein Nasr, ``The Place of the School of Isfahan in Islamic Philosophy and Sufism,'' pp. 3-19, and Leonard Lewisohn, ``Sufism and the School of Isfahan: Tasawwuf and Irfan in Late Safavid Iran,'' pp. 63-135.

\hyperref[_ednref6]{ vi } David

Lavender, \emph{The Southwest} (New York, 1980), p. 128; Robert L. Casey, \emph{Journey to the High Southwest} (Seattle, 1985), pp. 272-273, 276; New Mexico, various authors, pp. 123-124; Mason Sutherland
\& Justin Locke, ``Adobe New Mexico,'' \emph{National Geographic

Magazine}, December, 1949, pp. 823-824.

\hyperref[_ednref7]{ vii } Marta

Weigle, \emph{Brothers of Light, Brothers of Blood, The Penitentes of the Southwest} (Santa Fe, New Mexico, 1976), p. 32, citing Julio Puyol,
\emph{Platica de Disciplinantes} (Madrid, 1927), p. 245.


\hyperref[_ednref8]{ viii } Roy

Mottahedeh, \emph{The Mantle of the Prophet} (New York, 1985), pp. 175-176.


\flrightit{Posted on 2016-09-20 by Michael McClain}

\begin{center}* * *\end{center}

\begin{footnotesize}
\begin{sffamily}
\texttt{James K. Polk on 2016-09-27 at 01:54 said:}

Superb! Your website continues to contain some of the most substantive writings on the Internet. Keep up the good work and I will gladly continue reading!


\hfill

\end{sffamily}
\end{footnotesize}

